\section{Gaussian Space --- Invariance Distribution}
%%% generated by the blueprint pipeline from TCSlib.BooleanAnalysis.GaussianSpace.InvarianceDistribution
\subsection*{Gaussian Space --- Invariance Distribution: Introduction}

Declarations of the Lean module \texttt{TCSlib.BooleanAnalysis.GaussianSpace.InvarianceDistribution}.

\subsection*{Gaussian Space --- Invariance Distribution: Results}

\begin{lemma}[CDF bounds imply a Lévy distance bound]
\lean{BooleanAnalysis.GaussianSpace.Invariance.levyDistance_le_of_cdf}
\label{BooleanAnalysis.GaussianSpace.Invariance.levyDistance_le_of_cdf}
\leanok
\uses{BooleanAnalysis.GaussianSpace.levyDistance}
\difficulty{1}
For measures $\mu$ and $\nu$ on $\bbr$, define their Lévy distance as the infimum of all
$\delta\ge0$ such that, for every $u\in\bbr$,
$(\nu((-\infty,u-\delta])\.toReal-\delta)\le\mu((-\infty,u])\.toReal$ and
$\mu((-\infty,u])\.toReal\le\nu((-\infty,u+\delta])\.toReal+\delta$. If $\epsilon\ge0$
satisfies these two inequalities with $\delta=\epsilon$ for every $u$, then the Lévy
distance from $\mu$ to $\nu$ is at most $\epsilon$.
\end{lemma}

\begin{lemma}[Two-sided CDF comparison bounds Lévy distance]
\lean{BooleanAnalysis.GaussianSpace.Invariance.levyDistance_le_of_cdf_comparison}
\label{BooleanAnalysis.GaussianSpace.Invariance.levyDistance_le_of_cdf_comparison}
\leanok
\uses{BooleanAnalysis.GaussianSpace.Invariance.levyDistance_le_of_cdf, BooleanAnalysis.GaussianSpace.levyDistance}
\difficulty{4}
For finite measures $\mu$ and $\nu$ on $\bbr$, define their Lévy distance as the infimum
of all $\delta\ge0$ such that, for every $u\in\bbr$,
$(\nu((-\infty,u-\delta])\.toReal-\delta)\le\mu((-\infty,u])\.toReal$ and
$\mu((-\infty,u])\.toReal\le\nu((-\infty,u+\delta])\.toReal+\delta$. Suppose $w,e\ge0$
and, for every real $u$, each of the two measures' CDF values at $u$ is at most the
other's CDF value at $u+w$ plus $e$. Then the Lévy distance from $\mu$ to $\nu$ is at
most $w+e$.
\end{lemma}

\begin{lemma}[Integral bounds for a cutoff function]
\lean{BooleanAnalysis.GaussianSpace.Invariance.integral_cutoff_sandwich}
\label{BooleanAnalysis.GaussianSpace.Invariance.integral_cutoff_sandwich}
\leanok
\difficulty{5}
Let $\mu$ be a finite measure on $\mathbb R$, and let $\phi:\mathbb R\to\mathbb R$ be
continuous with $0\le\phi(x)\le1$ for every $x$. Suppose there are real numbers $u,v$
such that $\phi(x)=1$ whenever $x\le u$, and $\phi(x)=0$ whenever $v\le x$. Then
$\mu((-\infty,u])\le\int\phi\,d\mu\le\mu((-\infty,v])$, where the measure terms are
viewed as real numbers.
\end{lemma}

\begin{theorem}[Lévy distance bound from smooth test functions]
\lean{BooleanAnalysis.GaussianSpace.Invariance.levy_bound_at_scale}
\label{BooleanAnalysis.GaussianSpace.Invariance.levy_bound_at_scale}
\leanok
\uses{BooleanAnalysis.GaussianSpace.Invariance.exists_contDiff_cutoff, BooleanAnalysis.GaussianSpace.Invariance.integral_cutoff_sandwich, BooleanAnalysis.GaussianSpace.Invariance.levyDistance_le_of_cdf_comparison, BooleanAnalysis.GaussianSpace.levyDistance}
\difficulty{4}
Let $m$ be a nonnegative integer and let $\mu,\nu$ be probability measures on $\bbr$.
Define their Lévy distance as the infimum of all $\delta\ge0$ such that, for every
$u\in\bbr$, $(\nu((-\infty,u-\delta])\.toReal-\delta)\le\mu((-\infty,u])\.toReal$ and
$\mu((-\infty,u])\.toReal\le\nu((-\infty,u+\delta])\.toReal+\delta$. Suppose $B\ge0$
and, for every $(m+1)$ times continuously differentiable function $\phi:\bbr\to\bbr$,
every $C\ge0$ bounding $|\phi^{(m+1)}(t)|$ for all $t$, the difference of its integrals
against $\mu$ and $\nu$ has absolute value at most $CB$. Then for every $\eta>0$, the
Lévy distance is at most $(m+2)\eta+(2^m/\eta^{m+1})B$.
\end{theorem}

\begin{theorem}[Optimized scale bound]
\lean{BooleanAnalysis.GaussianSpace.Invariance.levy_bound_of_scale}
\label{BooleanAnalysis.GaussianSpace.Invariance.levy_bound_of_scale}
\leanok
\difficulty{6}
Let $m$ be a nonnegative integer and let $L,B\in\bbr$ with $B\ge0$. If, for every
$\eta>0$, $L\le(m+2)\eta+(2^m/\eta^{m+1})B$, then $L\le((m+2)+2^m)B^{1/(m+2)}$.
\end{theorem}

\begin{theorem}[Optimized Lévy distance bound]
\lean{BooleanAnalysis.GaussianSpace.Invariance.levy_bound_of_derivative}
\label{BooleanAnalysis.GaussianSpace.Invariance.levy_bound_of_derivative}
\leanok
\uses{BooleanAnalysis.GaussianSpace.Invariance.levy_bound_at_scale, BooleanAnalysis.GaussianSpace.Invariance.levy_bound_of_scale, BooleanAnalysis.GaussianSpace.levyDistance}
\difficulty{1}
Let $m$ be a nonnegative integer and let $\mu,\nu$ be probability measures on $\bbr$.
Define their Lévy distance as the infimum of all $\delta\ge0$ such that, for every
$u\in\bbr$, $(\nu((-\infty,u-\delta])\.toReal-\delta)\le\mu((-\infty,u])\.toReal$ and
$\mu((-\infty,u])\.toReal\le\nu((-\infty,u+\delta])\.toReal+\delta$. Suppose $B\ge0$
and, for every $(m+1)$ times continuously differentiable function $\phi:\bbr\to\bbr$,
every $C\ge0$ bounding $|\phi^{(m+1)}(t)|$ for all $t$, the absolute difference between
the integrals of $\phi$ against $\mu$ and $\nu$ is at most $CB$. Then the Lévy distance
is at most $((m+2)+2^m)B^{1/(m+2)}$.
\end{theorem}

\begin{theorem}[Lévy bound for mapped random variables]
\lean{BooleanAnalysis.GaussianSpace.Invariance.levy_bound_of_derivative_map}
\label{BooleanAnalysis.GaussianSpace.Invariance.levy_bound_of_derivative_map}
\leanok
\uses{BooleanAnalysis.GaussianSpace.Invariance.levy_bound_of_derivative, BooleanAnalysis.GaussianSpace.levyDistance}
\difficulty{3}
Let $m$ be a nonnegative integer, let $(\Omega,\mu)$ and $(\Xi,\nu)$ be probability
spaces, and let $U:\Omega\to\bbr$ and $V:\Xi\to\bbr$ be measurable almost everywhere
with respect to their respective measures. Suppose $B\ge0$ and, for every $(m+1)$ times
continuously differentiable function $\phi:\bbr\to\bbr$, every $C\ge0$ bounding
$|\phi^{(m+1)}(t)|$ for all $t$, the absolute difference between the integrals of
$\phi(U)$ with respect to $\mu$ and $\phi(V)$ with respect to $\nu$ is at most $CB$.
Then the Lévy distance between the distributions of $U$ and $V$ is at most
$((m+2)+2^m)B^{1/(m+2)}$.
\end{theorem}

\begin{lemma}[Fifth-root power inequality]
\lean{BooleanAnalysis.GaussianSpace.Invariance.invariance_fifth_power_le}
\label{BooleanAnalysis.GaussianSpace.Invariance.invariance_fifth_power_le}
\leanok
\uses{BooleanAnalysis.GaussianSpace.degree_mul_nine_pow_le_sixteen_pow}
\difficulty{4}
For every natural number $k$ and real number $\varepsilon\ge0$,
\[
\left(\frac{k\,9^k\varepsilon}{12}\right)^{1/5}\le 2^k\varepsilon^{1/5}.
\]
\end{lemma}

\begin{lemma}[Fourth-root power inequality]
\lean{BooleanAnalysis.GaussianSpace.Invariance.invariance_c3_quarter_power_le}
\label{BooleanAnalysis.GaussianSpace.Invariance.invariance_c3_quarter_power_le}
\leanok
\difficulty{5}
For every natural number $k$ and real numbers $r,\varepsilon$ with $6/5\le r$ and
$\varepsilon\ge0$,
\[
\left(\frac{k r^{3k}\sqrt{\varepsilon}}{3}\right)^{1/4}\le r^{2k}\varepsilon^{1/8}.
\]
\end{lemma}
